\begin{tabularx}{\textwidth}{L{36mm}X X}
\toprule
\textbf{Kontextuselem} & \textbf{Tartalma} & \textbf{Miért szükséges?} \\
\midrule
Azonosítók & Processzazonosító, szülőprocessz, tulajdonos, jogosultsági adatok. & Ezek alapján tudja az operációs rendszer megkülönböztetni és ellenőrizni a folyamatokat. \\
Állapot & Például új, futásra kész, futó, blokkolt, felfüggesztett vagy befejezett állapot. & Az ütemező ebből látja, hogy a folyamat kaphat-e CPU-t, vagy valamilyen eseményre vár. \\
Programsámláló és regiszterek & A következő utasítás címe, általános és speciális CPU-regiszterek, veremmutató. & Kontextusváltáskor ezeket menteni és visszaállítani kell, hogy a folyamat ott folytatódjon, ahol félbeszakadt. \\
Ütemezési adatok & Prioritás, CPU-használat, időszelet, várakozási idő, ütemezési sorokhoz tartozó hivatkozások. & Ezek alapján dönt az ütemező a futtatási sorrendről és a CPU-idő elosztásáról. \\
Memória-információ & Címtér, lap- vagy szegmenstábla adatai, memóriahatárok, jogosultságok. & A folyamat izolációjához, címleképzéséhez és védelméhez szükséges. \\
I/O és fájlinformáció & Nyitott fájlok, eszközök, I/O kérések, kommunikációs kapcsolatok. & A folyamat erőforrásait és függő műveleteit tartja nyilván. \\
Nyilvántartási adatok & Elhasznált CPU-idő, indulási idő, kvóták, statisztikák. & Teljesítményméréshez, elszámoláshoz és adminisztrációhoz használható. \\
\bottomrule
\end{tabularx}
